Gallstone status has been studied in relation to demographic, metabolic, laboratory, and body-composition characteristics \parencite{lammert2016gallstones,esen2024early}. Combining these measurements presents a supervised classification problem in which predictive discrimination must be considered alongside missed cases, false positives, and model interpretability. A small clinical dataset also raises a methodological question: does additional model complexity provide a consistent advantage over a regularised linear baseline?

This study evaluates that question using the published UCI Gallstone dataset. The outcome is gallstones present at assessment; the analysis does not model time to disease onset. Its purpose is to compare candidate classifiers and investigate the robustness of their internal performance, rather than establish a deployable diagnostic system.

The study makes three contributions within this dataset: a comparison of five predictive model families under shared, training-specific preprocessing and RFE--CART feature selection; LOPO internal validation yielding one held-out prediction per record; and retrospective analysis of errors recurring across model families. The comparison distinguishes probability discrimination from sensitivity and specificity at a common decision threshold.
